\sloppy
\section{Study Design}
\label{sec:study_design} 

\subsection{Research Questions}

This study aims to answer the following research questions:

\textbf{RQ1: How does the performance of fine-tuning LLMs compare to existing approaches on three unit testing tasks?} 

\textbf{RQ2: What is the impact of various factors (including model series, model architecture, and model size) on the performance of LLMs?}

\textbf{RQ3: How do fine-tuning approaches perform compared to prompt engineering approaches?
}


\subsection{Evaluation Metrics}

Following previous studies~\cite{watson_learning_2020, yu_automated_2022, tufano_unit_2021, hu_identify_2023}, we focus on two types of metrics: (1) runtime-based metrics, which focus on the correctness and behavior of generated code within execution environments, primarily for test generation tasks; and (2) text-based metrics, which assess the syntactic and semantic accuracy of generated outputs, used for assertion generation and test evolution tasks.

\subsubsection{Runtime-based Metrics.}
We leverage five metrics to evaluate LLMs' performance on the test generation task: syntax error, compilation error, failing test, passing test, and correct test.


\begin{itemize}[leftmargin=1em]
    \item \textbf{Syntax Error.} The test has syntax errors.
    
    \item \textbf{Build Error.} The test has correct syntax but fails to build.
    
    \item \textbf{Failing Test.} The test builds but fails due to wrong assertions or expected behavior.
    
    \item \textbf{Passing Test.} The test builds and passes.
    
    \item \textbf{Correct Test.} The test passes and covers the correct focal method.
    
\end{itemize}

\subsubsection{Text-based Metrics.}

We then leverage three widely used textual metrics to evaluate the performance of LLMs on the assertion generation and test evolution tasks.


\begin{itemize}[leftmargin=1em]
    \item \textbf{Exact Match Accuracy (EM).} The EM metric~\cite{hou_large_2023}, also known as perfect accuracy, measures the proportion of correctly predicted outputs that exactly match the ground truth. A prediction is considered correct only if each token in the output sequence matches the ground truth precisely.
    
    \item \textbf{BLEU.} BLEU~\cite{papineni_bleu_2002} is an automatic metric used to evaluate the syntactic similarity between the predicted sequence and the reference sequence. It measures n-gram overlap between the translation and references, incorporating precision and brevity to ensure appropriate length.
    
    \item \textbf{CodeBLEU.} CodeBLEU~\cite{ren_codebleu_2020} is a code-aware variant of BLEU specifically designed for evaluating code synthesis tasks. Unlike BLEU, CodeBLEU incorporates syntactic similarity through Abstract Syntax Tree (AST) information and semantic similarity through data-flow analysis. This makes CodeBLEU particularly effective for assessing the quality of code generation tasks.

\end{itemize}

\subsection{Selected Models}
\label{sec:selected_models}

In our selection of LLMs, we evaluate 37 state-of-the-art models based on their architectural complexity and performance on key benchmarks. 
These models fall into three major architectural categories: encoder-only, encoder-decoder, and decoder-only. 
Our final selection includes 20 distinct model series (e.g., CodeT5, CodeGen), comprising a total of 37 models, each varying in size and architecture. 
The sizes of these models range from 60 million parameters (e.g., CodeT5-small) to 16 billion parameters (e.g., CodeGen-16b-multi), allowing us to evaluate models under different computational constraints while balancing efficiency and task performance.
We provide a detailed list of our selected models as shown in Table~\ref{tab:model_selection}.

\begin{table}[htbp]
    \centering
    \small
    \caption{Selected models in this work.}
    \label{tab:model_selection}
    \resizebox{\textwidth}{!}{
    \begin{tabular}{ll}
        \toprule
        \textbf{Architecture} & \textbf{Models} \\
        \midrule
        \textbf{Encoder-only} & CodeBERT, GraphCodeBERT, UniXcoder\\
        \midrule
        \multirow{3}{*}{\textbf{Encoder-decoder}} & CodeT5 Series (CodeT5-small, CodeT5-base, CodeT5-large) \\
        & CodeT5+ Series (CodeT5+ 220m, CodeT5+ 770m) \\
        & PLBART Series (PLBART-base, PLBART-large) \\
        \midrule
        \multirow{10}{*}{\textbf{Decoder-only}} & CodeGPT, SantaCoder, DeciCoder, CodeShell, InCoder, StarCoder \\
        & Phi Series (Phi-1, Phi-2) \\
        & CodeGen Series (CodeGen-350m-multi, CodeGen-2b-multi, CodeGen-6b-multi, CodeGen-16b-multi) \\
        & CodeGen2 Series (CodeGen2-1b-p, CodeGen2-3.7b-p) \\
        & StarCoderBase Series (StarCoderBase-1b, StarCoderBase-3b, StarCoderBase-7b, StarCoderBase-15.5b) \\
        & StarCoder2 Series (StarCoder2-3b, StarCoder2-7b, StarCoder2-15b) \\
        & CodeLlama Series (CodeLlama-7b, CodeLlama-13b) \\
        & CodeGemma Series (CodeGemma-2b, CodeGemma-7b) \\
        & DeepSeek-Coder Series (DeepSeek-Coder-1.3b-base, DeepSeek-Coder-6.7b-base) \\
        \bottomrule
    \end{tabular}
    }
\end{table}


For decoder-only LLMs, we represent tasks as auto-regressive generation tasks. 
Given an initial prompt or source sequence, the decoder model predicts the next token by sampling from the probability distribution  $P(t_i|t_1, …, t_{i-1})$, where  $t_1, ..., t_{i-1} $ represent all the tokens generated up to position $i$. This process continues iteratively until the model generates an end-of-sequence token or reaches a predefined length. 
For encoder-decoder LLMs, we represent tasks as sequence-to-sequence translation tasks. 
In these models, fine-tuning is performed on a mapped source-target pair, denoted as $mst_i = \{s_i, t_i\}$, where $s_i$ is the source sequence and $t_i$ is the target sequence. 
The fine-tuning process learns the mapping $s_i \rightarrow t_i$ as a conditional probability $P(t_i | s_i)$.
For encoder-only LLMs, we manually extend them into encoder-decoder architectures by integrating a transformer decoder stack and treat them as encoder-decoder LLMs.

After training the LLMs for each task, we apply the beam search strategy when provided with a source sequence. 
This strategy generates the predicted target sequence by selecting words from multiple candidate sequences based on the probability distribution over the vocabulary. 









    


    

    


    
    
    


    







    


\subsection{Dataset}
\label{sec:dataset}

\subsubsection{Test generation task.}
In the test generation task, we utilize {Methods2Test} dataset proposed by Tufano et al.~\cite{tufano_methods2test_2022, tufano_unit_2021} to fine-tune LLMs. {Methods2Test} is a comprehensive, supervised dataset comprising test cases paired with their corresponding focal methods. This dataset is created through an extensive mining process, resulting in 780,944 pairs of JUnit tests and focal methods.

However, due to resource constraints, we are unable to use the entire dataset. To reduce the dataset and enhance its quality, we filter the dataset by following rules.
(1) \textit{Length of Tokens.} To ensure data completeness and simplicity, we filter out entries with token lengths that are either too short or too long. Specifically, we retain entries where the length of the input focal method is between 64 and 2048, and the length of the output test case is between 16 and 512.  
(2) \textit{Construction.} To normalize the construction of test cases, we only retain entries that begin with the prefix \texttt{`@Test public void'} and do not throw exceptions.
(3) \textit{Duplicated Test Cases.} We find half of the focal methods in the dataset include at least two related test cases. To increase the variety of focal methods, we randomly select one pair from each set of test cases associated with the same focal methods.
(4) \textit{Filtered repository.} To maintain the quality of the dataset, we filter out any repository with fewer than 50 pairs. Conversely, to ensure a diverse set of test cases that are not dominated by a few repositories, we randomly sample 200 pairs from any repository with more than 200 pairs.

Applying the above rules to the dataset, we obtain 56,132 pairs of data and call this dataset $Method2Test_{filter}$. We then split the filtered dataset by repositories into training, validation, and test sets using an 8:1:1 ratio for further fine-tuning.


Furthermore, previous works~\cite{tufano_unit_2021, alagarsamy_a3test_2023, ni_casmodatest_2024} primarily adopt runtime-based metrics to evaluate the performance of test generation approaches, focusing on the correct test rate.
Similar to \textsc{AthenaTest}~\cite{tufano_unit_2021}, we employ Defects4J as our benchmark dataset to evaluate LLMs' performance and use the same five evaluated projects: Apache Common Cli~\cite{2023CLI}, Apache Common Csv~\cite{2023CSV}, Google Gson~\cite{2023gson}, JFreeChart~\cite{2023Chart} and Apache Commons Lang~\cite{2023Lang}. 
These projects cover various domains, such as command-line interfaces, data processing, serialization, visualization, and utilities. 



\subsubsection{Assertion generation task.}
In the assertion generation task, we use the dataset known as $Data_{old}$. This dataset is derived from a raw dataset used by \textit{ATLAS}~\cite{watson_learning_2020}. Each entry in $Data_{old}$ is referred to as a Test-Assert Pair (TAP). A TAP consists of two components: (1) a focal-test pair, which includes a test method without an assertion and its corresponding focal method; (2) assertions. 
To simplify the problem, $Data_{old}$ excludes any TAP where the assertions contain tokens that are not present in the focal-test pair. 
In total, $Data_{old}$ contains 156,760 TAPs, which are divided into training, validation, and test sets by the ratio of 8:1:1 for further fine-tuning.


\subsubsection{Test evolution task.}
In the test evolution task, we utilize the dataset proposed by Hu et al.~\cite{hu_identify_2023} Specifically for the obsolete test updating task mentioned in the paper. Each sample in the dataset consists of $<$\textit{original method, updated method, original test, updated test}$>$. The dataset contains 5,196 samples, which the authors originally split into training and test sets using a 9:1 ratio. We further create a validation set by randomly selecting 11\% of the samples from the training set, resulting in a final split ratio of 8:1:1 for the training, validation and test sets. 

























\subsection{Implementation Details}
\label{sec:implementation}



We implement all model training approaches using the PyTorch framework. For the studied models, we use the Hugging Face versions. For model training, we adopt the default training parameters from previous studies for each task. However, due to the variety of models and tasks, we specifically adjust some training parameters. We use the AdamW optimizer with a learning rate of 5e-5 and train for a maximum of 50 epochs for all tasks. Additionally, we apply early stopping with a patience of two epochs, meaning training halts if the loss does not decrease for two consecutive epochs.
Maximum input and output sequence lengths are adjusted according to task-specific and model-specific requirements. For the assertion generation and test evolution tasks, we set source and target sequence lengths to 512 and 256, respectively. For the test generation task, we use lengths of 2048 and 512 due to the long context information. 

We conduct our experiments on a computing cluster with several A100 nodes (8x NVIDIA A100-PCIE-40GB). We utilize DeepSpeed to accelerate training and reduce memory usage.





